\section{方法优势与机理解释}

\subsection{口径一致性带来的“评测可控性”}

本文将观测算子 \( H \) 与训练退化算子 \( \mathrm{DC} \) 绑定为同一实现与同一参数镜像，本质上把“评测口径”显式纳入训练闭环：模型不再只追求在数学真值域 \( (|\hat{u}-u|) \) 上变好，而是同时被约束在**观测域** \( (|H(\tilde{u})-y|) \) 上变好，从而减少训练与评测之间的隐性域偏差（domain shift）。这一做法与近年来对 PINN/物理学习数值稳定性与误差来源的系统分析所强调的“训练目标必须与实际约束一致”的观点一致。\cite{Cuomo2022}

\subsection{三件套损失的“互补约束结构”}

\begin{itemize}
    \item \(\mathbf{L_{\text{rec}}}\)（z-score 域点误差）提供局部逐点一致性，保证整体拟合能力；
    \item \(\mathbf{L_{\text{spec}}}\)（低频谱一致性）在大尺度结构上施加约束，缓解仅优化点误差时常见的谱域偏置与结构性漂移；
    \item \(\mathbf{L_{\text{dc}}}\)（原值域观测一致性）将训练目标锚定到评测口径，使 \( H_{\text{err}} \equiv ||H(\hat{y}) - y|| \) 与 Rel-L2 的变化趋势更一致，从而降低“指标断裂”。
\end{itemize}

三者形成“真值域—谱域—观测域”的互补约束，能解释第6章中观察到的同步下降现象（当 H/DC 复用成立时尤为明显）。

\subsection{统一接口与确定性训练提高“横向可比性”}

统一接口（输入打包与 `forward` 契约）+ 固定切分/标准化 + 配置快照/环境指纹，使得不同模型（FNO/DeepONet/ConvLSTM/Transformer 变体）之间的差异更集中地反映为“方法差异”，而不是工程差异。这一点与科学机器学习领域对可复现协议与误差归因的主流建议相一致。\cite{Cuomo2022}

\section{局限性与风险点（更具体的“何时会变差”）}

\subsection{复杂几何与非周期边界下的谱域瓶颈}

当边界条件强非周期、几何复杂或存在非结构网格时，谱域算子（或频域模块）的归纳偏置可能与真实物理/边界不匹配，容易在边界带出现误差积累与伪影扩散。此时需要更强的边界处理（边界带掩码、显式边界编码、几何嵌入、域分解）或采用对不规则域更友好的算子设计。相关方向已有针对“域/边界泛化”的神经算子变体讨论，可作为后续增强路径。

\subsection{极端稀疏与高噪声下的权重敏感性}

当观测极端稀疏或噪声显著时，\( L_{\text{dc}} \) 会把噪声也当作“必须拟合的观测约束”，若 \( \lambda_{\text{dc}} \) 设定不当，可能造成：
\begin{itemize}
    \item 观测域过拟合（\( H_{\text{err}} \) 看似下降但真值域结构退化）；
    \item 训练不稳定（梯度被噪声主导）。
\end{itemize}

因此需要引入鲁棒一致性策略（例如观测域的鲁棒损失、噪声模型、或不确定性加权）。Bayesian/PINN 不确定性量化方向为“噪声—泛化”折衷提供了可行框架。

\subsection{资源与流程成本上升}

三件套损失（尤其 FFT 与一致性退化）+ 资源四项统计 + 显著性检验，会带来额外计算与工程成本。若应用场景更偏工程部署，需要在“统计完备性”与“训练吞吐”之间做分级：例如保留最小必要的口径一致性与关键指标，把全面显著性检验作为定期离线评估。

\section{适用条件与工程实践建议（可操作）}

\begin{itemize}
    \item 先锁口径，再调模型：先固定 \( H/\mathrm{DC} \) 与一致性脚本通过，再做网络结构与损失权重的比较；否则横向对比没有可解释性基础。
    \item \( \lambda_{\text{dc}} \) 与噪声水平联动：噪声越大，越需要降低 \( \lambda_{\text{dc}} \) 或使用鲁棒观测一致性；否则 \( H_{\text{err}} \) 可能“好看但不可靠”。
    \item 低频阈值扫描应与任务尺度匹配：若下游更关注大尺度结构（如涡核/波包），可适当提高低频权重；若需要精细边界层/尖峰结构，需同步强化高频恢复策略（位置编码、分层训练、或多尺度损失）。
    \item 复杂边界优先做边界带诊断：对非周期/复杂几何任务，建议在评测中固定输出边界带放大图与边界带 RMSE（bRMSE），作为是否需要“几何/边界增强”的触发器。
\end{itemize}

\section{未来工作（与本文框架强耦合的方向）}

\subsection{主动学习与自适应采样（信息效率）}

在稀疏观测场景下，与其固定采样，不如利用误差/不确定性指导采样位置与时间点，形成“采样—重建—再采样”的闭环。残差/误差驱动的自适应采样研究已显示其对 PINN 训练效率与稳定性具有显著影响，可迁移到本文的观测一致性框架中。

\subsection{弱式/变分物理约束与复杂边界融合}

针对复杂几何与边界，弱式/变分形式的物理约束（如 VPINNs 思路）能在一定程度上降低对点式残差的敏感性，并提供更可控的数值性质；与本文的 \( L_{\text{dc}} \) 可形成互补：一个约束物理一致性，一个约束观测口径一致性。\cite{Berrone2022}

\subsection{数据基准与跨任务扩展}

PDEBench 提供了跨 PDE、跨初边界条件的统一基准，有利于系统评估“口径一致性 + 三件套损失”是否具备跨任务稳健性。建议后续扩展到更多 PDE 类别与更极端的观测退化，以形成更强的外推结论。\cite{takamoto2022pdebench}

\section{小结}

本文方法的关键价值不在于“某个网络结构更强”，而在于把**评测口径**变成可控变量：通过 H/DC 复用与观测一致性损失，将训练目标与评测目标对齐；再以低频谱一致性损失稳定大尺度结构恢复，并用统一接口与确定性训练保证横向可比与可复现。局限性主要集中在复杂边界、强噪声与资源成本；对应的增强路径包括不确定性加权、弱式物理约束与主动采样闭环。

\section{参考文献（APA）}

\begin{itemize}
    \item Cuomo, S., Di Cola, V. S., Giampaolo, F., Rozza, G., Raissi, M., \& Piccialli, F. (2022). Scientific machine learning through physics-informed neural networks: Where we are and what’s next. \textit{Journal/Conference survey (Acta Numerica related review materials)}. doi:10.1017/S0962492923000089 \url{https://doi.org/10.1017/S0962492923000089}
    \item Takamoto, M., Praditia, T., Leiteritz, R., MacKinlay, D., Alesiani, F., Pflüger, D., \& Niepert, M. (2022). \textit{PDEBench: An extensive benchmark for scientific machine learning}. arXiv:2210.07182. \url{https://arxiv.org/abs/2210.07182}
    \item Wang, S., Yu, X., \& Perdikaris, P. (2022). \textit{When and why PINNs fail to train: A neural tangent kernel perspective}. \textit{Journal of Computational Physics, 449}, 110768. \url{https://doi.org/10.1016/j.jcp.2021.110768}
    \item Wu, C., Zhu, M., Tan, Q., Kartha, Y., \& Lu, L. (2022). \textit{A comprehensive study of non-adaptive and residual-based adaptive sampling for physics-informed neural networks}. arXiv:2207.10289.
    \item Berrone, S., et al. (2022). \textit{Variational physics-informed neural networks (VPINNs)}. \textit{Journal of Scientific Computing}. \url{https://doi.org/10.1007/s10915-022-01950-4}
    \item De Hoop, M. V., et al. (2023). \textit{Domain-agnostic Fourier neural operators}. (Irregular domains/boundary handling).
    \item Linka, K., et al. (2024). Efficient Bayesian physics-informed neural networks: Screening, backpropagation, and uncertainty quantification. \textit{Journal of Computational Physics}.
\end{itemize}
